% GROK REVIEW::GROK-0712::2026-10-14T07:12:00+09:00
\documentclass[11pt]{article}
\usepackage[margin=1in]{geometry}
\usepackage{amsmath,amssymb}
\title{Independent Grok Review --- 2026-10-14 07:12 KST}
\author{Grok}
\date{14 October 2026}

\begin{document}
\maketitle

\begin{center}
\fbox{\parbox{0.9\linewidth}{
\textbf{Executive status: P versus NP remains OPEN.}
No load-bearing implication is closed by this review.
}}
\end{center}

\section{Scope}
This independent review examines the living document at repository-root \texttt{main.tex} (TWOTREEBRIDGE1 whole-document reconciliation, 2026-10-14), the specialist reports, the two-tree compatibility bridge, the shared-DAG cofactor construction, and the prior affine-residual materials. It confirms that all current firewalls remain correctly classified and that the unrestricted whole-computation lower bound stays OPEN. Status is unchanged in its load-bearing claims since the 06:12 KST review; the new bridge is retained as CANDIDATE pending independent specialist review.

\section{Confirmation of current spine}
The affine base, residual contraction, cotree witness bijection, and upper-bound escapes remain PROVED/EXACT. Prior cube selectors and TRACEISO1's deterministic root-isolation theorem remain PROVED/EXACT at their reviewed scopes. The one-variable shared-circuit elimination now has an explicit constant-factor cofactor construction with charged physical allocation and reverse witnesses; repeating it yields only an upper bound, not unavoidable growth. The new complete source-level reduction from graph 3-colouring to two-tree compatibility (TWOTREEBRIDGE1) is retained as CANDIDATE pending independent specialist review. Both trees admit linear-size boundary summaries; their conjunction is small; the residual cost is joint existential search over a total encoding of length \(9+h^2+2h\). This is a complete polynomial reduction, not a proof that the problem lies outside P. Growing-factor amplification remains CONDITIONAL. The unrestricted whole-computation lower bound remains OPEN.

No contradiction with prior exact results is found. The two-tree construction correctly relocates known NP-complete compatibility onto two explicitly given trees with small summary circuits; it is not hardness amplification. Representation-driven structure continues to serve as a useful firewall rather than a hardness source for the whole machine. The conditional audit under the assumption P=NP correctly shows that no retained restricted-model lower bound is contradicted by a polynomial-size circuit for partial-coloring extension queries.

\section{Open load-bearing arrows}
The implication from an arbitrary eligible short program (or uniform decider) to superpolynomial cost on a structured residual object remains OPEN. No model-matched all-algorithm theorem exists. The firewalls relativize and algebrize; no nonrelativizing hinge is known. The two-tree compatibility problem, while NP-complete under the stated reduction, does not by itself establish a superpolynomial lower bound for unrestricted computation. Opaque representations, arbitrary joint nonlinear compatibility, and historical compiler interfaces remain outside the current selector theorems without additional bridges.

\section{Highest-value next target}
The single highest-value next assignment remains appropriate: audit repeated shared-DAG elimination on the explicit two-tree reductions, charging every cofactor support, alternative source-aware recomputation and witness lift before proposing an unavoidable-work invariant. Bridge reviews both witness directions and the total encoding; Tests checks independent target solutions; Redteam attacks the circuit-size/search distinction; Lower Bound must justify any claimed universal cost independently of the chosen compiler. Failure of one elimination rule is not a lower bound for other algorithms. A successful selector or elimination is a useful firewall; it does not close a whole-machine implication.

\section{Status and recommendations}
P versus NP remains open. The living document is updated to record this independent review. All claims continue to be classified under the allowed statuses. The comments folder receives this review artifact.

% END_OF_VERIFIED_REVIEW::GROK::GROK-0712
\end{document}
